% Ethical considerations (optional section, outside the page count) and the anonymized
% Reproducibility paragraph. Mirrors docs/CHECKLIST_DRAFT.md and docs/DESIGN_DECISIONS.md 11.
% Unnumbered: it carries NO \label (a \label after \section* resolves to the preceding numbered
% section) and is referred to by name ("the Ethical Considerations section").
% STATUS: protocol. Nothing described here has happened yet; the closing \placeholder forces a
% re-read against what was actually done.
\section*{Ethical Considerations}

\paragraph{Human participants.}
Three native-speaker validators and about twenty human-baseline respondents take part. All are adults who volunteer, read a bilingual (Vietnamese and English) consent form describing the task, its duration and their right to withdraw at any time, and are not paid; we follow the precedent of community-built resources whose contributors were volunteers \citep{2024-acl-long-620,2020-findings-emnlp-195} and say so plainly. Participants are recruited by personal invitation from the authors' community and stand in no dependency relation to the authors. We collect no names, e-mail addresses or free text that could identify a person with the judgments; validators are identified by a letter whose key is held separately and deleted after publication, and forms record only a region tag and coarse demographics (age band, region where the participant grew up, years lived in Vietnam, whether raised abroad) for the data statement \citep{Q18-1041}. No participant information is ever sent to an API. The study is not reviewed by an institutional review board, because the authors have access to none; we answer the corresponding checklist item ``No'' with this justification, and we judge the risk to participants minimal: the task is a language judgment on wordplay items, participation is anonymous, every item with a known offensive reading is excluded from the human forms, and validators are warned that some items may have vulgar readings. Participants under 18 are not recruited, so no parental-permission procedure is needed.

\paragraph{Offensive content and dual use.}
Nói lái is a vehicle for satire and for hiding taboo meanings, and a benchmark of it necessarily contains items whose output, or whose decoding, is vulgar. A blocklist of taboo syllables and syllable pairs, seeded by the authors and extended by the validators, is applied at generation to gold answers, \task{T2} readings and \task{T3} candidates, per syllable and per syllable pair in either order, because ordinary inputs can have obscene swaps, and to an input only when the input pair itself is a taboo phrase, since several blocklisted syllables are ordinary words in context; attested items are screened by hand. Flagged items are excluded from the public development split, from every prompt sent to a hosted model and from the human forms; they stay, flagged, in the gated test file so that scoring is complete, and the validators' offensive-reading judgments on the sample give the residual rate. A model that is fluent in nói lái could, in principle, slip taboo content past a keyword filter; the generator is a few hundred lines of rules any speaker could write, the same engine is the detector (decode both swaps and check the blocklist), and gated release is the mitigation.

\paragraph{Data and terms of use.}
The spelling list and the word list consulted at generation time are GPL-licensed; the release states which of their entries the lexical items reproduce and under which terms (\placeholder{written permission from their authors, or GPLv2 with attribution}). XCOPA is CC~BY~4.0 and is used under that license; its re-encoded variants are produced at run time and never stored. The development split is released under CC~BY~4.0 and the test split is gated, encrypted and released under CC~BY-NC-ND~4.0 with a canary string; the per-item score files are public, while raw model outputs, which quote test items, are released only inside the gated repository under the same terms. Three hosted providers' free tiers are used; all require users to be 18 or older. A provider whose tier uses submitted content for training never receives test items: it runs either on a paid key with a data-use opt-out or on a development-derived set of the same shape, reported separately and never in Table~\ref{tab:main}, and the core goes only to providers whose terms exclude training on inputs, recorded per provider in the run manifests. All API experiments run on the account of an eligible adult, which the run manifests record by role; at most the core items are ever sent, and the sealed split never is. We use the providers' services to evaluate their models, which their terms permit, and use no output to train a model.

\paragraph{Generative assistance.}
Items and gold answers in \noilai{} come from a rule engine, not from a language model, and no metric uses an LLM judge. Code in the released repository, its documentation and drafts of this manuscript are produced with an AI coding assistant under the authors' direction, then run, tested and reviewed by the authors, who are responsible for every line and every claim; the research questions, design, rule tables and analyses are the authors'. The use is disclosed in the Responsible NLP checklist and will be stated in the Acknowledgements of the final version, so that the disclosure does not compromise anonymity now. Every reference is checked by hand against its source before submission.

\paragraph{Reproducibility.}
The generator, parser, re-encoder, prompt templates, scoring code, statistics and the probing and patching code are released at \anonurl, together with the development split, the gated test split, the attested seed table, the tokenizer-audit results and the per-item score files (item ids, covariates and scores); raw model outputs quote gated test items and are released only inside the gated test repository. Every dataset build and every model run writes a manifest with item counts and drop counts per cell (the dataset's build seed is withheld, because one seeded stream draws and splits the development and test items; its content hash is its identity), resource hashes, checkpoint identifier and revision, tokenizer and chat-template hashes, serving engine and version, requested and resolved precision, quantization, engine flags, prompt-file hashes, canary check, normalization census, hardware, wall time and GPU-hours; the manifests are released with the outputs. Compute totals \placeholder{TODO} GPU-hours on Tesla T4s and \placeholder{TODO} TPU-hours on a TPU v5e-8, all on free tiers, plus \placeholder{TODO} API calls; the log is released. Randomness is seeded everywhere; a build is byte-identical for a given seed and resource set. The hypotheses, metrics, exclusion rules, stopping rules and analysis plan are pre-registered before any test-split run in a dated commit whose hash is given in Appendix~\ref{app:prereg}, and every deviation from it is logged with a date.

\placeholder{Re-read every paragraph above against what was actually done before submission: who took part and how they were recruited, which providers and accounts were used, the compute totals, the reference check and the pre-registration commit; every present-tense protocol sentence becomes a past-tense report or is corrected.}
